\chapter{KESIMPULAN DAN SARAN}

\section{Kesimpulan}
Berdasarkan hasil implementasi, pengujian, dan evaluasi terhadap penambahan dukungan protokol FINS pada sistem HMI FUXA, dapat disimpulkan hal-hal berikut ini sesuai dengan rumusan masalah penelitian:

\begin{enumerate}
    \item \textbf{Penambahan fitur protokol FINS pada sistem HMI FUXA.}  
    Pengembangan modul konektor FINS berhasil dilakukan dengan menambahkan komponen \texttt{DeviceFins} pada sisi \textit{backend} (Node.js) dan antarmuka konfigurasi pada sisi \textit{frontend} (Angular).  
    Modul ini mampu menjalankan fungsi utama komunikasi PLC, yaitu \texttt{connect()}, \texttt{read()}, \texttt{write()}, dan \texttt{poll()} dengan protokol FINS berbasis UDP maupun TCP.  
    Implementasi ini memungkinkan FUXA berkomunikasi langsung dengan PLC Omron CJ2M tanpa memerlukan middleware tambahan, menjadikannya solusi HMI \textit{open source} pertama yang mendukung FINS secara native.

    \item \textbf{Perbaikan proses pembacaan dan penulisan data FINS agar stabil dan dapat divisualisasikan pada UI.}  
    Berdasarkan hasil uji fungsional dan beban, sistem mampu mempertahankan tingkat keberhasilan komunikasi sebesar 99\% dengan latensi rata-rata baca/tulis sebesar 48\,ms.  
    Nilai ini menunjukkan performa komunikasi yang stabil untuk penggunaan laboratorium atau sistem produksi berskala kecil–menengah.  
    Data hasil pembacaan ditampilkan secara real-time di antarmuka HMI melalui mekanisme event-driven FUXA (\texttt{device-value:changed} dan \texttt{device-status:changed}), yang memastikan sinkronisasi antara nilai tag PLC dan visualisasi pengguna.  
    Selain itu, fitur auto-reconnect dan logging adaptif yang ditambahkan meningkatkan keandalan komunikasi ketika terjadi gangguan jaringan.

\end{enumerate}

Secara keseluruhan, penelitian ini membuktikan bahwa FUXA dapat dikembangkan menjadi sistem HMI yang kompatibel dengan PLC Omron menggunakan protokol FINS.  
Integrasi ini berhasil meningkatkan interoperabilitas sistem tanpa menambah biaya lisensi perangkat lunak.  
Namun, hasil pengujian menunjukkan bahwa deteksi alarm masih memiliki latensi sekitar 1000\,ms, lebih lambat dari target 200\,ms, sehingga diperlukan optimasi polling atau mekanisme notifikasi berbasis event dari PLC untuk memenuhi kebutuhan real-time yang lebih ketat.

\section{Saran}
Berdasarkan hasil dan keterbatasan penelitian, beberapa saran pengembangan dapat diajukan untuk meningkatkan kinerja dan penerapan sistem di masa mendatang:

\begin{enumerate}
    \item \textbf{Optimasi deteksi alarm dan latensi real-time.}  
    Implementasikan pendekatan berbasis event (\textit{event-driven listener}) untuk tag alarm kritis atau gunakan fitur \textit{interrupt task} dari PLC Omron agar pembacaan tidak bergantung pada interval polling reguler.

    \item \textbf{Pengujian lanjutan dan tuning performa.}  
    Lakukan \textit{stress test} dengan jumlah tag lebih besar dan interval polling lebih cepat untuk menentukan batas performa sistem.  
    Data hasil pengujian dapat digunakan untuk menyempurnakan algoritma scheduler FUXA, mengurangi beban CPU, dan mencegah penundaan pada UI.

    \item \textbf{Penguatan aspek keamanan dan akses jaringan.}  
    Tambahkan fitur autentikasi pengguna, enkripsi TLS/DTLS, dan opsi VPN agar komunikasi antara HMI dan PLC lebih aman, terutama untuk implementasi di jaringan produksi.

    \item \textbf{Integrasi penyimpanan historis dan analitik.}  
    Integrasikan basis data deret waktu (\textit{time-series database}) seperti InfluxDB atau TimescaleDB untuk menyimpan data tag secara historis, mendukung visualisasi tren, alarm historis, dan analisis performa mesin.

    \item \textbf{Peningkatan antarmuka pengguna dan pengalaman operator.}  
    Tambahkan fitur seperti tombol \textit{Test Connection}, validasi parameter otomatis (SA1, DA1, alamat memori), serta umpan balik status koneksi agar sistem lebih ramah bagi pengguna non-teknis.

    \item \textbf{Ekspansi kompatibilitas protokol.}  
    Lanjutkan pengembangan agar konektor FINS juga mendukung seri PLC Omron lainnya (mis. NX/NJ) serta varian komunikasi FINS over TCP.  
    Selain itu, perlu diteliti integrasi multi-protokol (Modbus, MQTT, OPC-UA) untuk memperluas interoperabilitas industri.

    \item \textbf{Publikasi kode sumber dan dokumentasi.}  
    Dokumentasikan seluruh modul, skrip pengujian, serta dataset hasil uji dalam repositori \textit{open source} agar penelitian dapat direplikasi dan dikembangkan lebih lanjut oleh komunitas industri dan akademik.
\end{enumerate}

\noindent
Dengan penerapan saran-saran di atas, FUXA berpotensi menjadi sistem HMI/SCADA sumber terbuka yang tidak hanya mendukung PLC Omron melalui protokol FINS, tetapi juga memenuhi kebutuhan performa, keamanan, dan skalabilitas di lingkungan industri nyata.
